\subsection{\label{sec.details.gf.laure}Laurent power series}

\begin{fineprint}
Next, we shall prove some claims made in Section \ref{sec.gf.laure}. We shall
not go into too many details here, since the proofs are mostly analogous to
the corresponding proofs in the theory of \textquotedblleft
usual\textquotedblright\ FPSs (with nonnegative exponents), which we have
already seen. However, every once in a while, some minor change is required;
we will mainly focus on these changes.

\begin{proof}
[Proof of Theorem \ref{thm.fps.laure.laupol-ring} (sketched).]First, we need
to prove that the set $K\left[  x^{\pm}\right]  $ is closed under addition,
under scaling, and under the multiplication introduced in Definition
\ref{def.fps.laure.laupol}. The proof of this is analogous to Theorem
\ref{thm.fps.pol.ring} (with the obvious changes, such as replacing the
$\sum_{i=0}^{n}$ sums by $\sum_{i\in\mathbb{Z}}$ sums).

Next, we need to prove that the multiplication on $K\left[  x^{\pm}\right]  $
is associative, commutative, distributive and $K$-bilinear. This is analogous
to the corresponding parts of Theorem \ref{thm.fps.ring}. Likewise, we can
show that the element $\left(  \delta_{i,0}\right)  _{i\in\mathbb{Z}}$ (which
is our analogue of the FPS $\underline{1}$) is a neutral element for this
multiplication. Hence, $K\left[  x^{\pm}\right]  $ is a commutative
$K$-algebra with unity $\left(  \delta_{i,0}\right)  _{i\in\mathbb{Z}}$.

It remains to show that the element $x$ is invertible in this $K$-algebra. But
this is easy: Set $\overline{x}:=\left(  \delta_{i,-1}\right)  _{i\in
\mathbb{Z}}$, and show (by direct calculation) that $x\overline{x}%
=\overline{x}x=1$. (For example, let us prove that $x\overline{x}=1$. Indeed,
from $x=\left(  \delta_{i,1}\right)  _{i\in\mathbb{Z}}$ and $\overline
{x}=\left(  \delta_{i,-1}\right)  _{i\in\mathbb{Z}}$, we obtain%
\[
x\overline{x}=\left(  \delta_{i,1}\right)  _{i\in\mathbb{Z}}\cdot\left(
\delta_{i,-1}\right)  _{i\in\mathbb{Z}}=\left(  c_{n}\right)  _{n\in
\mathbb{Z}},\ \ \ \ \ \ \ \ \ \ \text{where}\ \ \ \ \ \ \ \ \ \ c_{n}%
=\sum_{i\in\mathbb{Z}}\delta_{i,1}\delta_{n-i,-1}%
\]
(by Definition \ref{def.fps.laure.laupol}). However, for each $n\in\mathbb{Z}%
$, we have%
\begin{align*}
c_{n}  &  =\sum_{i\in\mathbb{Z}}\delta_{i,1}\delta_{n-i,-1}=\underbrace{\delta
_{1,1}}_{=1}\delta_{n-1,-1}+\sum_{\substack{i\in\mathbb{Z};\\i\neq
1}}\underbrace{\delta_{i,1}}_{\substack{=0\\\text{(since }i\neq1\text{)}%
}}\delta_{n-i,-1}\\
&  \ \ \ \ \ \ \ \ \ \ \ \ \ \ \ \ \ \ \ \ \left(
\begin{array}
[c]{c}%
\text{here, we have split off the addend for }i=1\\
\text{from the sum}%
\end{array}
\right) \\
&  =\delta_{n-1,-1}+\underbrace{\sum_{\substack{i\in\mathbb{Z};\\i\neq
1}}0\delta_{n-i,-1}}_{=0}=\delta_{n-1,-1}=\delta_{n,0}%
\ \ \ \ \ \ \ \ \ \ \left(
\begin{array}
[c]{c}%
\text{since }n-1=-1\text{ holds}\\
\text{if and only if }n=0
\end{array}
\right)  .
\end{align*}
Thus, $\left(  c_{n}\right)  _{n\in\mathbb{Z}}=\left(  \delta_{n,0}\right)
_{n\in\mathbb{Z}}=1$, so that $x\overline{x}=\left(  c_{n}\right)
_{n\in\mathbb{Z}}=1$, as we wanted to prove. The proof of $\overline{x}x=1$ is
analogous, or follows from $x\overline{x}=1$ by commutativity.)
\end{proof}
\end{fineprint}

\begin{fineprint}
Proposition \ref{prop.fps.laure.a=sumaixi} is an analogue of Corollary
\ref{cor.fps.sumakxk}. Recall that we proved the latter corollary using Lemma
\ref{lem.fps.xa} and Proposition \ref{prop.fps.xk}. Thus, in order to prove
the former proposition, we need to find analogues of Lemma \ref{lem.fps.xa}
and Proposition \ref{prop.fps.xk} for Laurent polynomials.

The analogue of Lemma \ref{lem.fps.xa} for Laurent polynomials is the following:

\begin{lemma}
\label{lem.fps.laure.xa}Let $\mathbf{a}=\left(  a_{n}\right)  _{n\in
\mathbb{Z}}$ be a Laurent polynomial in $K\left[  x^{\pm}\right]  $. Then,
\[
x\cdot\mathbf{a}=\left(  a_{n-1}\right)  _{n\in\mathbb{Z}}%
\ \ \ \ \ \ \ \ \ \ \text{and}\ \ \ \ \ \ \ \ \ \ x^{-1}\cdot\mathbf{a}%
=\left(  a_{n+1}\right)  _{n\in\mathbb{Z}}.
\]

\end{lemma}

\begin{proof}
[Proof of Lemma \ref{lem.fps.laure.xa} (sketched).]From $x=\left(
\delta_{i,1}\right)  _{i\in\mathbb{Z}}$ and $\mathbf{a}=\left(  a_{n}\right)
_{n\in\mathbb{Z}}=\left(  a_{i}\right)  _{i\in\mathbb{Z}}$, we obtain%
\[
x\cdot\mathbf{a}=\left(  \delta_{i,1}\right)  _{i\in\mathbb{Z}}\cdot\left(
a_{i}\right)  _{i\in\mathbb{Z}}=\left(  c_{n}\right)  _{n\in\mathbb{Z}%
},\ \ \ \ \ \ \ \ \ \ \text{where}\ \ \ \ \ \ \ \ \ \ c_{n}=\sum
_{i\in\mathbb{Z}}\delta_{i,1}a_{n-i}%
\]
(by Definition \ref{def.fps.laure.laupol}). However, for each $n\in\mathbb{Z}%
$, we have%
\begin{align*}
c_{n}  &  =\sum_{i\in\mathbb{Z}}\delta_{i,1}a_{n-i}=\underbrace{\delta_{1,1}%
}_{=1}a_{n-1}+\sum_{\substack{i\in\mathbb{Z};\\i\neq1}}\underbrace{\delta
_{i,1}}_{\substack{=0\\\text{(since }i\neq1\text{)}}}a_{n-i}%
\ \ \ \ \ \ \ \ \ \ \left(
\begin{array}
[c]{c}%
\text{here, we have split off the}\\
\text{addend for }i=1\\
\text{from the sum}%
\end{array}
\right) \\
&  =a_{n-1}+\underbrace{\sum_{\substack{i\in\mathbb{Z};\\i\neq1}}0a_{n-i}%
}_{=0}=a_{n-1}.
\end{align*}
Thus, $\left(  c_{n}\right)  _{n\in\mathbb{Z}}=\left(  a_{n-1}\right)
_{n\in\mathbb{Z}}$, so that $x\cdot\mathbf{a}=\left(  c_{n}\right)
_{n\in\mathbb{Z}}=\left(  a_{n-1}\right)  _{n\in\mathbb{Z}}$. Thus we have
proved%
\begin{equation}
x\cdot\mathbf{a}=\left(  a_{n-1}\right)  _{n\in\mathbb{Z}}.
\label{pf.lem.fps.laure.xa.1}%
\end{equation}


It remains to prove that $x^{-1}\cdot\mathbf{a}=\left(  a_{n+1}\right)
_{n\in\mathbb{Z}}$. Here we use a trick: Let $\mathbf{b}=\left(
a_{n+1}\right)  _{n\in\mathbb{Z}}$; this is again a Laurent polynomial in
$K\left[  x^{\pm}\right]  $. Hence, (\ref{pf.lem.fps.laure.xa.1}) (applied to
$\mathbf{b}$ and $a_{n+1}$ instead of $\mathbf{a}$ and $a_{n}$) yields%
\[
x\cdot\mathbf{b}=\left(  a_{\left(  n-1\right)  +1}\right)  _{n\in\mathbb{Z}%
}=\left(  a_{n}\right)  _{n\in\mathbb{Z}}%
\]
(since $a_{\left(  n-1\right)  +1}=a_{n}$ for each $n\in\mathbb{Z}$). In other
words, $x\cdot\mathbf{b}=\mathbf{a}$ (since $\mathbf{a}=\left(  a_{n}\right)
_{n\in\mathbb{Z}}$). Dividing this equality by the invertible element $x$, we
find $\mathbf{b}=x^{-1}\cdot\mathbf{a}$, so that $x^{-1}\cdot\mathbf{a}%
=\mathbf{b}=\left(  a_{n+1}\right)  _{n\in\mathbb{Z}}$. This completes the
proof of Lemma \ref{lem.fps.laure.xa}.
\end{proof}

The analogue of Proposition \ref{prop.fps.xk} for Laurent polynomials is the following:

\begin{proposition}
\label{prop.fps.laure.xk}We have%
\[
x^{k}=\left(  \delta_{i,k}\right)  _{i\in\mathbb{Z}}%
\ \ \ \ \ \ \ \ \ \ \text{for each }k\in\mathbb{Z}.
\]

\end{proposition}

\begin{proof}
[Proof of Proposition \ref{prop.fps.laure.xk} (sketched).]This is similar to
Proposition \ref{prop.fps.xk}, which we proved by induction on $k$. Here, too,
we can use induction, but (since $k$ can be negative) we have to use
\textquotedblleft two-sided induction\textquotedblright, which contains both
an induction step from $k$ to $k+1$ and an induction step from $k$ to $k-1$.
(See \cite[\S 2.15]{detnotes} for a detailed explanation of two-sided induction.)

Both induction steps rely on Lemma \ref{lem.fps.laure.xa}. (Specifically, the
step from $k$ to $k+1$ uses the $x\cdot\mathbf{a}=\left(  a_{n-1}\right)
_{n\in\mathbb{Z}}$ part of Lemma \ref{lem.fps.laure.xa}, whereas the step from
$k$ to $k-1$ uses the $x^{-1}\cdot\mathbf{a}=\left(  a_{n+1}\right)
_{n\in\mathbb{Z}}$ part.)
\end{proof}

Now, Proposition \ref{prop.fps.laure.a=sumaixi} is easy:

\begin{proof}
[Proof of Proposition \ref{prop.fps.laure.a=sumaixi} (sketched).]Analogous to
Corollary \ref{cor.fps.sumakxk}, but using Proposition \ref{prop.fps.laure.xk}
instead of Proposition \ref{prop.fps.xk}.
\end{proof}

Next, let us prove Theorem \ref{thm.fps.laure.lauser-ring}:

\begin{proof}
[Proof of Theorem \ref{thm.fps.laure.lauser-ring} (sketched).]This is
analogous to the proof of Theorem \ref{thm.fps.laure.laupol-ring}. The only
(slightly) different piece is the proof that the set $K\left(  \left(
x\right)  \right)  $ is closed under multiplication. So let us prove this:

Let $\left(  a_{n}\right)  _{n\in\mathbb{Z}}$ and $\left(  b_{n}\right)
_{n\in\mathbb{Z}}$ be two elements of $K\left(  \left(  x\right)  \right)  $.
We must prove that their product $\left(  a_{n}\right)  _{n\in\mathbb{Z}}%
\cdot\left(  b_{n}\right)  _{n\in\mathbb{Z}}$ belongs to $K\left(  \left(
x\right)  \right)  $ as well. This product is defined by%
\[
\left(  a_{n}\right)  _{n\in\mathbb{Z}}\cdot\left(  b_{n}\right)
_{n\in\mathbb{Z}}=\left(  c_{n}\right)  _{n\in\mathbb{Z}}%
,\ \ \ \ \ \ \ \ \ \ \text{where}\ \ \ \ \ \ \ \ \ \ c_{n}=\sum_{i\in
\mathbb{Z}}a_{i}b_{n-i}.
\]
Thus, we need to prove that $\left(  c_{n}\right)  _{n\in\mathbb{Z}}$ belongs
to $K\left(  \left(  x\right)  \right)  $. In other words, we must prove that
the sequence $\left(  c_{-1},c_{-2},c_{-3},\ldots\right)  $ is essentially
finite (by the definition of $K\left(  \left(  x\right)  \right)  $).

So let us prove this now. We know that the sequence $\left(  a_{-1}%
,a_{-2},a_{-3},\ldots\right)  $ is essentially finite (since $\left(
a_{n}\right)  _{n\in\mathbb{Z}}\in K\left(  \left(  x\right)  \right)  $).
Thus, there exists some negative integer $p$ such that%
\begin{equation}
\text{all }i\leq p\text{ satisfy }a_{i}=0.
\label{pf.thm.fps.laure.lauser-ring.3}%
\end{equation}
Similarly, there exists some negative integer $q$ such that%
\begin{equation}
\text{all }j\leq q\text{ satisfy }b_{j}=0
\label{pf.thm.fps.laure.lauser-ring.4}%
\end{equation}
(since $\left(  b_{n}\right)  _{n\in\mathbb{Z}}\in K\left(  \left(  x\right)
\right)  $). Consider these $p$ and $q$. Now, set $r:=p+q$. We shall show that
all negative integers $n\leq r$ satisfy $c_{n}=0$.

Indeed, let $n\leq r$ be any integer. Then, each integer $i\geq p$ satisfies
$\underbrace{n}_{\leq r=p+q}-\underbrace{i}_{\geq p}\leq p+q-p=q$ and
therefore
\begin{equation}
b_{n-i}=0 \label{pf.thm.fps.laure.lauser-ring.4b}%
\end{equation}
(by (\ref{pf.thm.fps.laure.lauser-ring.4}), applied to $j=n-i$). Now,%
\[
c_{n}=\sum_{i\in\mathbb{Z}}a_{i}b_{n-i}=\sum_{\substack{i\in\mathbb{Z}%
;\\i<p}}\underbrace{a_{i}}_{\substack{=0\\\text{(by
(\ref{pf.thm.fps.laure.lauser-ring.3}))}}}b_{n-i}+\sum_{\substack{i\in
\mathbb{Z};\\i\geq p}}a_{i}\underbrace{b_{n-i}}_{\substack{=0\\\text{(by
(\ref{pf.thm.fps.laure.lauser-ring.4b}))}}}=\underbrace{\sum_{\substack{i\in
\mathbb{Z};\\i<p}}0b_{n-i}}_{=0}+\underbrace{\sum_{\substack{i\in
\mathbb{Z};\\i\geq p}}a_{i}0}_{=0}=0.
\]


Forget that we fixed $n$. We thus have shown that all $n\leq r$ satisfy
$c_{n}=0$. Hence, the sequence $\left(  c_{-1},c_{-2},c_{-3},\ldots\right)  $
is essentially finite. As explained, this completes our proof of the claim
that the set $K\left(  \left(  x\right)  \right)  $ is closed under multiplication.

The rest of Theorem \ref{thm.fps.laure.lauser-ring} is proved just like
Theorem \ref{thm.fps.laure.laupol-ring}.
\end{proof}
\end{fineprint}
